% Part: sets-functions-relations
% Chapter: relations
% Section: orders

\documentclass[../../../include/open-logic-section]{subfiles}

\begin{document}

\olfileid{sfr}{rel}{ord}
\olsection{Ordeningen}

\begin{explain}
Als we objecten vergelijken, zeggen we vaak dat ze in een bepaald opzicht
``kleiner dan'', ``gelijk aan'' of ``groter dan'' andere objecten zijn. Daarvoor
gebruiken we \emph{ordeningsrelaties}. Er zijn verschillende soorten. Bij
sommige kun je elk paar objecten met elkaar vergelijken; bij andere kan dat
niet. Bij sommige telt gelijkheid mee (zoals bij~$\le$), bij andere niet (zoals
bij~$<$). Het helpt dus om deze relaties in soorten in te delen.
\end{explain}

\begin{defn}[Preorde]
Een relatie is een \emph{preorde} als ze reflexief én transitief is.  
\end{defn}

\begin{defn}[Partiële ordening]
Een preorde is een \emph{partiële ordening} als ze ook antisymmetrisch is.
\end{defn}

\begin{defn}[Lineaire ordening]\ollabel{def:linearorder}
Een partiële ordening is een \emph{totale ordening} of \emph{lineaire
ordening} als ze ook samenhangend is.
\end{defn}

\begin{ex}
Elke lineaire ordening is ook een partiële ordening. Elke partiële ordening is
weer een preorde. Maar andersom klopt geen van beide. De universele relatie
op~$A$ is een preorde, want ze is reflexief en transitief. Als $A$ meer dan één
!!{element} heeft, is deze relatie niet antisymmetrisch. Daarom is ze geen
partiële ordening.
\end{ex}

\begin{ex}
Neem de relatie \emph{niet langer dan} $\preccurlyeq$ op~$\Bin^*$:
$x \preccurlyeq y$ geldt precies wanneer $\len{x} \le \len{y}$. Deze relatie
is reflexief en transitief, dus een preorde. Ze is zelfs samenhangend. Toch is
ze geen partiële ordening, want ze is niet antisymmetrisch. Bijvoorbeeld: $01
\preccurlyeq 10$ en $10 \preccurlyeq 01$, maar $01 \neq 10$.
\end{ex}

\begin{ex}
Een belangrijk voorbeeld van een partiële ordening is de relatie $\subseteq$ op
een verzameling van verzamelingen. Dit is niet altijd een lineaire ordening.
Neem $a \neq b$ en bekijk
$\Pow{\{a, b\}} = \{\emptyset, \{a\}, \{b\}, \{a,b\}\}$. Dan geldt
$\{a\} \nsubseteq \{b\}$, $\{a\} \neq \{b\}$ en $\{b\}
\nsubseteq \{a\}$.
\end{ex}

\begin{ex}
De relatie \emph{deelbaarheid zonder rest} is een partiële ordening, maar geen
lineaire ordening. Voor gehele getallen $n$ en~$m$ betekent $n \mid m$: $n$
deelt $m$ zonder rest. Dat is precies het geval als er een geheel getal~$k$
bestaat waarvoor $m = kn$. Op~$\Nat$ is dit een partiële ordening, maar geen
lineaire ordening: zo geldt $2 \nmid 3$ en ook $3 \nmid 2$. Op $\Int$ is
deelbaarheid alleen een preorde, want de relatie is daar niet antisymmetrisch:
$1 \mid -1$ en $-1 \mid 1$, maar $1 \neq -1$.
\end{ex}

\begin{ex}
Op de verzameling rijen~$A^*$ gebruiken we de volgende
\emph{uitbreidingsrelatie}: $s \sqsubseteq s'$ geldt precies wanneer
$s = \emptyseq$ (de lege rij), $s = s'$, of $s =
\tuple{s_1, \dots, s_n}$ en $s' =
\tuple{s_1, \dots, s_n, s_{n+1}, \dots, s_m}$. Als $s \sqsubseteq s'$,
noemen we $s$ ook een \emph{beginsegment} van~$s'$. Deze relatie op $A^*$ is
een partiële ordening, maar geen lineaire ordening. Als bijvoorbeeld $a \neq b$,
dan geldt $ab \not\sqsubseteq ba$ en $ba \not\sqsubseteq ab$.
\end{ex}

\begin{defn}[Strikte ordening]
Een relatie is een \emph{strikte ordening} als ze irreflexief, asymmetrisch en
transitief is.
\end{defn}

\begin{defn}[Strikte lineaire ordening]\ollabel{def:strictlinearorder}
Een strikte ordening is een \emph{strikte totale ordening} of \emph{strikte
lineaire ordening} als ze ook samenhangend is.
\end{defn}

\begin{ex}
De lineaire ordening~$\le$ hoort bij de strikte lineaire ordening~$<$. De
partiële ordening~$\subseteq$ hoort bij de strikte ordening~$\subsetneq$.
\end{ex}

We kunnen van elke strikte ordening $R$ op~$A$ een partiële ordening maken.
Daarvoor voegen we de diagonaal $\Id{A}$ toe: alle paren~$\tuple{x,
x}$.  (Dit heet de \emph{reflexieve afsluiting} van~$R$.)
Andersom maken we van een partiële ordening een strikte ordening door
~$\Id{A}$ te verwijderen. De volgende twee resultaten leggen precies vast hoe
dit werkt.

\begin{prop}\ollabel{prop:stricttopartial}
Als $R$ een strikte ordening op~$A$ is, is $R^+ = R \cup \Id{A}$ een partiële
ordening. Als $R$ ook een strikte lineaire ordening is, is $R^+$ een lineaire
ordening.
\end{prop}

\begin{proof}
Neem aan dat $R$ een strikte ordening is. Dit betekent dat $R \subseteq A^2$ en
dat $R$ irreflexief, asymmetrisch en transitief is. Stel
$R^+ = R \cup \Id{A}$. We moeten bewijzen dat $R^+$ reflexief,
antisymmetrisch en transitief is.

$R^+$ is reflexief, want $\tuple{x, x} \in \Id{A} \subseteq R^+$ voor elke
$x \in A$. 

We bewijzen nu dat $R^+$ antisymmetrisch is. Neem, om een tegenspraak te krijgen,
aan dat $R^+xy$ en $R^+yx$, maar $x \neq y$. Er geldt
$\tuple{x,y} \in R \cup \Id{A}$, terwijl $\tuple{x, y} \notin \Id{A}$.
Daarom moet $\tuple{x, y} \in R$ gelden, dus $Rxy$. Op dezelfde manier volgt
$Ryx$. Dat is in strijd met de aanname dat $R$ asymmetrisch is.

Voor de transitiviteit nemen we aan dat $R^+xy$ en $R^+yz$. Als
$\tuple{x, y} \in R$ en $\tuple{y,z} \in R$, dan is $\tuple{x, z} \in R$,
want $R$ is transitief. Anders geldt $\tuple{x, y} \in \Id{A}$, dus $x = y$,
of $\tuple{y, z} \in \Id{A}$, dus $y = z$. In het eerste geval weten we al dat
$R^+yz$; samen met $x = y$ geeft dit $R^+xz$. In het tweede geval gaat het
net zo. In beide gevallen geldt $R^+xz$. Daarmee is $R^+$ ook transitief.

Neem nu ook aan dat $R$ samenhangend is.
Voor alle $x \neq y$ geldt dan $Rxy$ of~$Ryx$. Anders gezegd:
$\tuple{x, y} \in R$ of $\tuple{y, x} \in R$. Omdat $R \subseteq R^+$,
geldt dit ook voor $R^+$. Dus is ook $R^+$ samenhangend.
\end{proof}

\begin{prop}\ollabel{prop:partialtostrict}
Als $R$ een partiële ordening op~$A$ is, is $R^- = R \setminus \Id{A}$ een
strikte ordening. Als $R$ ook een lineaire ordening is, is $R^-$ een strikte
lineaire ordening.
\end{prop}

\begin{proof}
Dit bewijs laten we als opgave.
\end{proof}

\begin{prob}
Bewijs \olref[sfr][rel][ord]{prop:partialtostrict}. 
\end{prob}

Het volgende eenvoudige resultaat zegt dat strikte lineaire ordeningen een
eigenschap hebben die op extensionaliteit lijkt:

\begin{prop}\ollabel{prop:extensionality-strictlinearorders}
Als $<$ een strikte lineaire ordening op $A$ is, dan geldt:
\[
  (\forall a, b \in A)((\forall x \in A)(x < a \liff x < b) \lif a = b).
\]
\end{prop}

\begin{proof}
Neem aan dat $(\forall x \in A)(x < a \liff x < b)$. Als $a < b$, dan zou
$a < a$ gelden. Dat kan niet, want $<$ is irreflexief. Dus $a \nless b$.
Om precies dezelfde reden geldt $b \nless a$. Dan moet $a = b$, want $<$ is
samenhangend.
\end{proof}
\end{document}
